\section*{Capítulo \ref{ch:b3:structural-biology} --- Biología estructural de las proteínas}

\begin{solution}{exo:b3:structural-biology:1}
Primaria: la secuencia de los $141$ residuos de una cadena $\alpha$.
Secundaria: sus ocho hélices~$\alpha$. Terciaria: el \omterm{def:b3:structural-biology:levels}{plegamiento} de
globina de una cadena envolviendo su hemo. Cuaternaria: el tetrámero
$\alpha_{2}\beta_{2}$ y la manera en que sus dímeros giran uno contra
otro.
\end{solution}

\begin{solution}{exo:b3:structural-biology:2}
$36\times \qty{0.15}{nm} = \qty{5.4}{nm}$; $36/3.6 = 10$ vueltas;
$36 - 4 = 32$ puentes de hidrógeno.
\end{solution}

\begin{solution}{exo:b3:structural-biology:3}
Que la estructura nativa, incluido el emparejamiento correcto de cuatro
puentes disulfuro de entre $105$ posibilidades, se recupera
espontáneamente a partir de la cadena desplegada: la secuencia contiene
toda la información necesaria y el estado nativo es el mínimo
termodinámico. Hacerlo en un tubo de ensayo con proteína pura excluía
cualquier molde, enzima o maquinaria celular como fuente de la
información.
\end{solution}

\begin{solution}{exo:b3:structural-biology:4}
RMN: proteínas pequeñas, por debajo de unos \qty{40}{kDa}, en disolución.
Cristalografía: cualquier tamaño que cristalice, de los péptidos a los
ribosomas. \omterm{def:b3:structural-biology:methods}{Criomicroscopia electrónica}: complejos grandes, por encima de
unos \qty{50}{kDa}, sin cristales. La RMN informa directamente del
movimiento (y la crio-ME puede ordenar las moléculas en clases
conformacionales); una estructura cristalina es un promedio estático.
\end{solution}

\begin{solution}{exo:b3:structural-biology:5}
$2^{150} \approx 1.4\times 10^{45}$ conformaciones, $1.4\times 10^{33}$
s $\approx 5\times 10^{25}$ años. $3^{20} \approx 3.5\times 10^{9}$,
$3.5$ ms: el péptido podría buscar de forma exhaustiva en un segundo; la
proteína no podría en la edad del universo.
\end{solution}

\begin{solution}{exo:b3:structural-biology:6}
$Y = [\alpha(1+\alpha) + Lc\alpha(1+c\alpha)]/[(1+\alpha)^{2} +
L(1+c\alpha)^{2}]$. $\alpha = 1$: $3.01/106 = 0.028$ (un solo sitio,
$0.50$). $\alpha = 10$: $121/242 = 0.50$ ($0.91$). $\alpha = 100$:
$\num{10300}/\num{10601} = 0.97$ ($0.99$). La curva de MWC va muy por
detrás del sitio único con poco ligando y sube después bruscamente ---
de $0.03$ a $0.97$ a lo largo de dos décadas en las que el sitio único va
de $0.5$ a $0.99$: esa pendiente es la \omterm{def:b3:structural-biology:allostery}{cooperatividad}.
\end{solution}

\begin{solution}{exo:b3:structural-biology:7}
El bisfosfoglicerato se une solo al \omterm{def:b3:structural-biology:allostery}{estado T} (en la cavidad central), de
modo que lo estabiliza: en el modelo sube $L$, la razón $T/R$ de la
proteína sin ligando, y desplaza así la curva a la derecha --- menos
afinidad y más oxígeno liberado a las presiones de los tejidos. Los
glóbulos rojos almacenados pierden el compuesto, $L$ baja, la curva se
desplaza a la izquierda, y la sangre transfundida retiene su oxígeno en
lugar de cederlo en los tejidos, hasta que las células lo resintetizan a
lo largo de un día.
\end{solution}

\begin{solution}{exo:b3:structural-biology:8}
$d = \lambda/(2\sin\theta) = 0.1/(2\times 0.259) = \qty{0.19}{nm}$: el
esqueleto y todas las cadenas laterales quedan situados, con los átomos
individuales resueltos. Para \qty{0.12}{nm}:
$\sin\theta = 0.1/0.24 = 0.417$ y $\theta =
\qty{25}{\degree}$.
\end{solution}

\begin{solution}{exo:b3:structural-biology:9}
Una lisina cargada enterrada entre cadenas laterales apolares cuesta
decenas de kilojulios por mol (su carga sin solvatar, perdido el
enterramiento hidrófobo de la leucina), más que todo el $\Delta G$ del
\omterm{def:b3:structural-biology:levels}{plegamiento}: la proteína se despliega o el núcleo se reorganiza. En la
superficie la lisina queda hidratada igual que lo estaría en el estado
desplegado y no se pierde nada.
$\qty{-30}{kJ/mol} \approx -12RT$ significa $K = e^{12} \approx 1.6\times
10^{5}$: una molécula de cada \num{160000} está desplegada en cada
instante, y todo el margen equivale a dos o tres puentes de hidrógeno de
entre centenares --- lo que vale una sola mutación.
\end{solution}

\begin{solution}{exo:b3:structural-biology:10}
Las especies $R$ suman $[R_{0}](1+\alpha)^{n}$ y las $T$,
$[R_{0}]L(1+c\alpha)^{n}$, de modo que $\bar R = (1+\alpha)^{n}/[(1+\alpha)^{n}
+ L(1+c\alpha)^{n}]$. Para $c \to 0$ el término de $T$ es $L$, y $\bar R =
1/2$ cuando $(1+\alpha)^{n} = L$, es decir, $\alpha = L^{1/n} - 1$. Para
la hemoglobina $(3\times 10^{5})^{1/4} = 23.4$, luego
$\alpha \approx 22$: \qty{22}{mmHg}, cerca del $P_{50}$ de
\qty{26}{mmHg} --- el conmutador conformacional y la semisaturación casi
coinciden.
\end{solution}

\begin{solution}{exo:b3:structural-biology:11}
Se resistió porque todo agente infeccioso conocido se replicaba mediante
ácido nucleico, y una proteína no puede copiar su secuencia. Encontrar un
ácido nucleico necesario, o mostrar que la proteína purificada no
transmite, la refutaría. Apoyos: la infectividad sobrevive a las
nucleasas y a la radiación ultravioleta e ionizante a dosis que destruyen
cualquier genoma, pero la destruyen los desnaturalizantes de proteínas y
las proteasas; la PrP recombinante plegada in vitro en fibrillas
transmite la enfermedad a animales, con las mismas propiedades de cepa al
pasarla. Hace falta una semilla porque la conversión de un monómero
normal es prohibitivamente lenta por sí sola --- primero ha de formarse
el núcleo de la estructura $\beta$ cruzada, que es el suceso limitante
---, mientras que el extremo de una fibrilla ya formada convierte
monómeros deprisa: la forma se copia por crecimiento sobre molde.
\end{solution}

\begin{solution}{exo:b3:structural-biology:12}
Muchas secuencias se pliegan en la misma estructura: lo que un
\omterm{def:b3:structural-biology:levels}{plegamiento} necesita es un patrón de posiciones hidrófobas y polares,
unas pocas glicinas y prolinas en los giros y los residuos que hacen el
trabajo; todo lo demás de la superficie puede cambiar sin penalización,
de modo que las mutaciones se acumulan ahí mientras la selección
conserva el \omterm{def:b3:structural-biology:levels}{plegamiento} y la función. Los residuos conservados son por
tanto los del núcleo (más como patrón que como identidades), los
catalíticos y de unión, y las glicinas, prolinas y cisteínas
estructurales. Los métodos de perfil ponderan justamente esas columnas y
reconocen así miembros de la familia al \qty{15}{\%} de identidad; el
diseño de proteínas recorre la lógica al revés, eligiendo un patrón
hidrófobo para un \omterm{def:b3:structural-biology:levels}{plegamiento} diana y dejando que la superficie sea
cualquier cosa.
\end{solution}

\begin{solution}{pb:b3:structural-biology:1}
\textbf{1.} $0.75\times 150 \approx 112$ residuos helicoidales; $112\times
\qty{0.15}{nm} = \qty{17}{nm}$; $112/3.6 \approx 31$ vueltas.
\textbf{2.} $112 - 8\times 4 = 80$ puentes de hidrógeno.
\textbf{3.} Masa $150\times 110 = \qty{16500}{Da} = 2.7\times 10^{-20}$
g; volumen $2.7\times 10^{-20}\times 0.73 = 2.0\times 10^{-20}$ cm$^{3}$
$= \qty{20}{nm^3}$; radio $(3V/4\pi)^{1/3} = \qty{1.7}{nm}$.
\textbf{4.} Extendida, $150\times 0.36 = \qty{54}{nm}$ frente a un
diámetro de \qty{3.4}{nm}: un factor de $16$.
\textbf{5.} Los residuos helicoidales son el \qty{75}{\%} de la longitud
extendida, $112\times 0.36 = \qty{40}{nm}$; la hélice los comprime a
\qty{17}{nm}, un factor $0.36/0.15 = 2.4$.
\textbf{6.} $40\times 4 = \qty{160}{kJ/mol}$ ganados, $150\times 0.9 =
\qty{135}{kJ/mol}$ perdidos: en neto,
$\Delta G \approx \qty{-25}{kJ/mol}$, dentro del intervalo marginal.
\textbf{7.} $3^{150} \approx 4\times 10^{71}$; $4\times 10^{59}$ s
$\approx 10^{52}$ años.
\textbf{8.} $10^{-2}\times 10^{12} = 10^{10}$ conformaciones, una
fracción $3\times 10^{-62}$ del recuento.
\textbf{9.} Las hélices, en \qty{1}{\micro s} (en paralelo); los
$8! = \num{40320}$ empaquetamientos a $10^{6}$ por segundo tardan
\qty{40}{ms}: en total unos \qty{40}{ms}, el orden de magnitud
observado. Sí: resolver las partes por separado y luego el conjunto es un
embudo --- el espacio de búsqueda se factoriza, no se agota.
\textbf{10.} Ciclos esperados, $1/0.3 = 3.3$, unos \qty{33}{s}; tras seis
ciclos queda sin plegar $0.7^{6} = 0.12$.
\textbf{11.} $\Delta G = -RT\ln K$ con $K = 10^{4} \to 10^{2}$: de
$-23.7$ a \qty{-11.9}{kJ/mol}, una pérdida de \qty{12}{kJ/mol} de
estabilidad ($RT = \qty{2.58}{kJ/mol}$, $\ln 100 = 4.6$).
\textbf{12.} La agregación empieza con un núcleo, cuya velocidad de
formación sube como una potencia alta de la concentración de la especie
desplegada propensa a agregarse; cien veces más de esa especie eleva
varios órdenes de magnitud la probabilidad de que se forme una semilla
dentro de una vida, y una vez que existe una semilla el crecimiento es
rápido y se automoldea.
\textbf{13.} $\alpha = 100$: $Y = 1.054\times 10^{8}/1.089\times 10^{8} =
0.97$. $\alpha = 20$: numerador $185\,220 + 103\,680 = 288\,900$,
denominador $194\,481 + 622\,080 = 816\,561$: $Y = 0.35$.
\textbf{14.} $\bar R = 1.041\times 10^{8}/1.089\times 10^{8} = 0.96$ en
el pulmón; $194\,481/816\,561 = 0.24$ en el músculo.
\textbf{15.} $0.97 - 0.35 = 0.62$ de la capacidad descargada (dos tercios
de la carga). Un solo sitio: $100/126 - 20/46 = 0.79 - 0.43 = 0.36$.
\textbf{16.} Capacidad $150\times 1.34 = \qty{200}{mL/L}$; al músculo,
$0.62\times 200 \approx \qty{125}{mL/L}$. En reposo, $(0.97 - 0.78)\times
200 = \qty{38}{mL/L}$; $250/38 \approx \qty{6.5}{L/min}$ de flujo
sanguíneo --- el orden de un gasto cardíaco en reposo.
\textbf{17.} $L = 3\times 10^{6}$, $\alpha = 40$: numerador $2.76\times
10^{6} + 3.29\times 10^{6} = 6.05\times 10^{6}$, denominador $2.83\times
10^{6} + 11.5\times 10^{6} = 14.4\times 10^{6}$: $Y = 0.42$ en lugar de
$0.78$ --- mucho más oxígeno descargado en reposo, la adaptación de la
sangre a la altitud y del músculo al trabajo.
\textbf{18.} Fetal, $L = 3\times 10^{4}$, $\alpha = 30$: numerador
$893\,730 + 19\,770 = 913\,500$, denominador $923\,520 + 85\,680 =
1\,009\,200$: $Y = 0.91$. Materna a la misma presión: numerador
$893\,730 + 197\,730 = 1\,091\,460$, denominador $923\,520 + 856\,830
= 1\,780\,350$: $Y = 0.61$. A igual presión la sangre fetal retiene más
oxígeno, de modo que el oxígeno fluye de la madre al feto a través de la
placenta.
\textbf{19.} $(10^{5}\,\text{nm}/6\,\text{nm})^{3} = 4.6\times 10^{12}$
celdas; $1.9\times 10^{13}$ moléculas.
\textbf{20.} $d = 0.1/(2\sin 14.5^{\circ}) = \qty{0.20}{nm}$: sí, todas
las cadenas laterales y todos los átomos quedan situados.
\textbf{21.} $\sin\theta = 0.1/0.3$, $\theta = \qty{19.5}{\degree}$. En
un cristal desordenado las moléculas no están en registro exacto, las
ondas dispersadas a ángulo alto --- que codifican el detalle fino --- ya
no se suman en fase y las manchas se desvanecen a ángulo alto; solo
sobreviven las manchas de baja \omterm{def:b3:structural-biology:methods}{resolución}.
\textbf{22.} Reducir $d$ a la mitad exige el doble de señal y por tanto
cuatro veces más partículas: unas \num{40000} (más en la práctica, pues
entran otros errores).
\textbf{23.} Las proteínas de membrana, que rara vez cristalizan a partir
de disoluciones de detergente y pueden fotografiarse en un nanodisco
lipídico; y los complejos grandes y flexibles --- ribosomas,
espliceosomas, canales iónicos con sus reguladores ---, cuya
heterogeneidad impide la cristalización pero cuyas partículas sí pueden
ordenarse en clases conformacionales.
\textbf{24.} El hemo y el oxígeno unidos y su geometría exacta, el agua y
los iones ordenados, la segunda conformación (T frente a R) y la prueba
de que el \omterm{def:b3:structural-biology:levels}{plegamiento} es correcto --- una predicción es un único modelo
sin ligando y sin comprobación experimental.
\textbf{25.} Se descarga $0.62$ de la capacidad (un solo sitio, $0.36$);
el tiempo de Levinthal, unos $10^{52}$ años; y un mapa a \qty{0.20}{nm}.
\end{solution}
